\section{Shortest Paths with Dijkstra's Algorithm}

Once the weighted visibility graph has been constructed, motion planning
becomes a shortest-path problem on a graph.

For each node $v$, Dijkstra's algorithm maintains a tentative distance

\[
d(v),
\]

representing the shortest currently known distance from the start node to
$v$.

Initially,

\[
d(S)=0,
\]

while every other node is assigned

\[
d(v)=\infty.
\]

The central operation in Dijkstra's algorithm is \emph{edge relaxation}.
For an edge from $u$ to $v$ with weight $w(u,v)$, a new candidate distance
to $v$ is

\[
d(u)+w(u,v).
\]

If

\[
d(u)+w(u,v)<d(v),
\]

then a shorter route to $v$ has been discovered, and we update

\[
d(v)=d(u)+w(u,v).
\]

The predecessor of $v$ is also recorded as $u$ so that the final shortest
path can later be reconstructed.

\subsection{Selecting the Next Node}

After relaxing the edges of the current node, Dijkstra's algorithm selects
the unvisited node with the smallest tentative distance.

Once a node is selected in this way, its shortest-path distance is considered
finalized.

This property relies on all graph edge weights being non-negative. In the
visibility graph, edge weights are Euclidean distances, so this condition is
automatically satisfied.

\subsection{Path Reconstruction}

Whenever relaxation improves the tentative distance to a node $v$, the
algorithm records the node $u$ from which that improvement came.

This node is called the \emph{predecessor} or \emph{parent} of $v$.

After the goal is reached, the shortest path is reconstructed by following
the parent links backward:

\[
G \rightarrow \operatorname{parent}(G)
\rightarrow \operatorname{parent}(\operatorname{parent}(G))
\rightarrow \cdots \rightarrow S.
\]

Because this produces the path in reverse order, the resulting sequence is
reversed to obtain the final path from start to goal.

\subsection{Undirected Visibility Edges}

The visibility graph is undirected: if node $u$ can see node $v$, then node
$v$ can also see node $u$.

For efficiency, each undirected edge is stored only once in the graph.

Therefore, when processing a node, Dijkstra's algorithm must check whether
the current node appears as either endpoint of an edge.